\chapter{Introduction}
\label{chap:introduction}

Credit-default prediction is an important decision-support task because an
incorrect decision can expose a lender to loss or restrict a reliable
borrower's access to credit. Ensemble models can learn non-linear relationships
in borrower data, but their decisions are harder to inspect than those of a
simple scorecard. A high aggregate score also does not show whether applicants
from different age groups receive comparable treatment.

This thesis studies prediction, explanation, and fairness as connected
properties of one credit-risk model. It uses LightGBM, SHAP explanations,
age-group fairness metrics, an age-blind model, and reweighing. The study also
audits data preparation because a preprocessing decision can create or hide a
fairness disparity before a model is trained.

\section{Motivation}

Credit-risk models must identify likely defaults while supporting review and
accountability. SHAP can show which features influence a prediction, but an
explanation does not automatically establish that the model treats groups
fairly. Conversely, an aggregate accuracy or ROC-AUC does not reveal group
differences. Removing age from the input also does not guarantee fairness when
credit history, employment, housing, or income carry age information.

\section{Problem Statement}

The problem addressed by this thesis is how to evaluate a credit-default model
beyond overall predictive performance. The study asks which features drive
the model, whether approvals differ between age groups, whether disparity
remains after age removal, and what performance cost is associated with
reweighing.

\section{Research Questions}
\label{sec:research-questions}

\begin{enumerate}
  \item Which features most influence the LightGBM model, and how do SHAP
  explanations compare with built-in gain and split importance?
  \item Does the model treat applicants under 25 and applicants aged 25 and
  above fairly under demographic parity and equal opportunity?
  \item Does age-related disparity remain after removing age, and which
  remaining features may act as age proxies?
  \item How much does reweighing reduce measured disparity, and what does it
  cost in accuracy and ROC-AUC?
\end{enumerate}

\section{Objectives}

\begin{enumerate}
  \item Develop and evaluate a LightGBM credit-default classifier on cleaned
  borrower data.
  \item Compare gain, split, and mean absolute SHAP importance.
  \item Measure demographic parity and equal opportunity with uncertainty
  intervals.
  \item Retrain without age and audit residual age information and candidate
  proxy features.
  \item Evaluate reweighing through both fairness changes and predictive-cost
  measurements.
\end{enumerate}

\section{Contributions}

The thesis contributes an integrated evaluation of predictive performance,
explainability, and age-group fairness. It documents a preparation audit that
found duplicate rows, age-selective outlier deletion, and engineered features
that recovered age exactly. It also defines a reproducible comparison between
an age-inclusive model, an age-blind model, and their reweighed variants.

\section{Scope and Limitations}

The study concerns offline prediction on the public Credit Risk Dataset. It
uses 11 recorded predictors, a binary default label, LightGBM, SHAP, the
under-25 versus 25-and-above grouping, and reweighing. The dataset is
simulated, its age distribution is narrow, and outcomes are available only for
granted loans. Therefore, the findings are specific to this dataset and
protocol and do not establish a general legal or fairness guarantee.

\section{Thesis Organization}

Chapter 2 reviews ensemble credit models, SHAP, group fairness, proxy
features, and mitigation. Chapter 3 presents the methodology and experimental
design. Chapter 4 describes the implementation and evaluation setup. Chapter
5 reports the available baseline and audit findings and defines the required
final result tables. Chapter 6 summarizes the contributions, objective status,
and future work.
